\newcommand{\DCR}{DeepCycleESN}
\newcommand{\DCRalt}{$\mathrm{DeepESN_{Cycle}$}}
\newcommand{\DAR}{DeepAntisymESN}
\newcommand{\DARalt}{$\mathrm{DeepESN_{Antisym}$}}
\newcommand{\DCO}{DeepCycleRON}
\newcommand{\DCOalt}{$\mathrm{DeepRON_{Cycle}$}}
\newcommand{\DAO}{DeepAntisymRON}
\newcommand{\DAOalt}{$\mathrm{DeepRON_{Antisym}$}}



\newcommand{\WtotCycle}{
    \begin{bmatrix}
    \Wrec{1} & \mathbf{0} & \cdots & \mathbf{\Wcyc{1}} \\
    \Wcyc{2} & \Wrec{2} & \cdots & \mathbf{0} \\
    \vdots & \ddots & \ddots & \vdots \\
    \mathbf{0} & \cdots & \Wcyc{L
    } & \Wrec{L}
    \end{bmatrix}
}

% add to the diagonal the W_rec^l 
\newcommand{\WtotAntisym}{
    \begin{bmatrix}
    \Wrec{1}-\gamma I & -\epsilon\Cbwd{1}^\top & \cdots & \mathbf{0} \\
    \epsilon\Cfwd{2} & \Wrec{2}-\gamma I & \ddots & \vdots \\
    \vdots & \ddots & \ddots & -\epsilon\Cbwd{L-1}^\top \\
    \mathbf{0} & \cdots & \epsilon\Cfwd{L} & \Wrec{L}-\gamma I
    \end{bmatrix}
}


\label{ch:appendix-A}

This appendix derives the Jacobian and necessary conditions for the Echo State Property (ESP) of the
\emph{Deep Cycle Echo State Network} (\DCR) described in Chapter~\ref{chap:methodology}. 

\section{Model definition}
\label{sec:dcr-definition}

Let the input be $\mathbf{u}(t)\in\mathbb{R}^{N_U}$ and consider $L$ reservoir modules with states
\[
\mathbf{h}^{(l)}(t)\in\mathbb{R}^{n_l},\qquad l=1,\dots,L.
\]
The global reservoir state introduced in Eq.~\eqref{eq:method:feature_concat} is the concatenation
\[
\mathbf{H}(t) \coloneqq \big[\mathbf{h}^{(1)}(t),\dots,\mathbf{h}^{(L)}(t)\big] \in \mathbb{R}^{N},\quad N=\sum_{l=1}^L n_l.
\]

All modules share the same leak rate $\alpha\in(0,1]$, and the element-wise nonlinearity is $\tanh(\cdot)$.
The pre-activation of the first module matches Eq.~\eqref{eq:method:cycle_dynamics}:
\begin{align}
\mathbf{z}^{(1)}(t) &=
\mathbf{W}_{\text{rec}}^{(1)} \mathbf{h}^{(1)}(t-1)
+ \mathbf{W}_{\text{in}}^{(1)} \mathbf{u}(t)
+ \mathbf{W}_{\text{cyc}}^{(L)} \mathbf{h}^{(L)}(t-1)
+ \mathbf{b}^{(1)}, \label{eq:dcr-z1}\\
\mathbf{z}^{(l)}(t) &=
\mathbf{W}_{\text{rec}}^{(l)} \mathbf{h}^{(l)}(t-1)
+ \mathbf{W}_{\text{in}}^{(l)} \mathbf{u}(t)
+ \mathbf{W}_{\text{cyc}}^{(l-1)} \mathbf{h}^{(l-1)}(t-1)
+ \mathbf{b}^{(l)},\qquad l=2,\dots,L. \label{eq:dcr-zl}
\end{align}
Each module is updated through the leaky integration rule
\begin{equation}
\mathbf{h}^{(l)}(t) = (1-\alpha)\,\mathbf{h}^{(l)}(t-1) + \alpha\,\tanh\big(\mathbf{z}^{(l)}(t)\big).
\label{eq:dcr-leak}
\end{equation}

\begin{remark} On Time indexing for the inter-module links.
All inter-module terms depend on $\mathbf{h}^{(\cdot)}(t-1)$, so the global map $\mathbf{H}(t)=F\big(\mathbf{u}(t),\mathbf{H}(t-1)\big)$ is well-defined without introducing algebraic loops, exactly as discussed in Section~\ref{sec:method:cycle}.
\end{remark}
\section{Jacobian of the global state map of  \DCR}
\label{sec:dcr-jacobian}

For a fixed input $\mathbf{u}(t)$, define the state transition map
\[
\mathbf{H}(t) = F\big(\mathbf{u}(t),\mathbf{H}(t-1)\big).
\]
The Jacobian with respect to the previous global state is
\[
\mathbf{J}(t) \coloneqq \frac{\partial \mathbf{H}(t)}{\partial \mathbf{H}(t-1)} \in \mathbb{R}^{N\times N}.
\]

\subsection{Layerwise derivatives}

Let
\[
\mathbf{D}^{(l)}(t) = \mathrm{diag}\!\big(\tanh'(\mathbf{z}^{(l)}(t))\big) \in \mathbb{R}^{n_l\times n_l}.
\]
From \eqref{eq:dcr-leak}--\eqref{eq:dcr-zl} we obtain, for the first module,
\begin{align}
\frac{\partial \mathbf{h}^{(1)}(t)}{\partial \mathbf{h}^{(1)}(t-1)}
&= (1-\alpha) \mathbf{I}_{n_1} + \alpha \mathbf{D}^{(1)}(t)\, \mathbf{W}_{\text{rec}}^{(1)}, \label{eq:dcr-j11}\\
\frac{\partial \mathbf{h}^{(1)}(t)}{\partial \mathbf{h}^{(L)}(t-1)}
&= \alpha \mathbf{D}^{(1)}(t)\, \mathbf{W}_{\text{cyc}}^{(L)}. \label{eq:dcr-j1L}
\end{align}
For every deeper module $l\ge 2$ we have
\begin{align}
\frac{\partial \mathbf{h}^{(l)}(t)}{\partial \mathbf{h}^{(l)}(t-1)}
&= (1-\alpha) \mathbf{I}_{n_l} + \alpha \mathbf{D}^{(l)}(t)\, \mathbf{W}_{\text{rec}}^{(l)}, \label{eq:dcr-jll}\\
\frac{\partial \mathbf{h}^{(l)}(t)}{\partial \mathbf{h}^{(l-1)}(t-1)}
&= \alpha \mathbf{D}^{(l)}(t)\, \mathbf{W}_{\text{cyc}}^{(l-1)}. \label{eq:dcr-jl-lminus1}
\end{align}

Introduce the block matrices
\begin{align}
\mathbf{R}^{(l)}(t) &\coloneqq (1-\alpha) \mathbf{I}_{n_l} + \alpha \mathbf{D}^{(l)}(t)\, \mathbf{W}_{\text{rec}}^{(l)}, \label{eq:dcr-R}\\
\mathbf{P}^{(l)}(t) &\coloneqq \alpha \mathbf{D}^{(l)}(t)\, \mathbf{W}_{\text{cyc}}^{(l-1)},\qquad l=2,\dots,L, \label{eq:dcr-P}\\
\mathbf{C}^{(1)}(t) &\coloneqq \alpha \mathbf{D}^{(1)}(t)\, \mathbf{W}_{\text{cyc}}^{(L)}. \label{eq:dcr-C}
\end{align}

\subsection{Block structure}

Let $\mathbf{J}_{l,k}(t)\coloneqq \partial \mathbf{h}^{(l)}(t)/\partial \mathbf{h}^{(k)}(t-1)$. Using \eqref{eq:dcr-j11}--\eqref{eq:dcr-jl-lminus1} we obtain
\begin{align}
\mathbf{J}_{1,1}(t) &= \mathbf{R}^{(1)}(t), &
\mathbf{J}_{1,L}(t) &= \mathbf{C}^{(1)}(t), &
\mathbf{J}_{1,k}(t) &= \mathbf{0}\quad (k\notin\{1,L\}), \label{eq:dcr-J-row1}\\
\mathbf{J}_{l,l}(t) &= \mathbf{R}^{(l)}(t), &
\mathbf{J}_{l,l-1}(t) &= \mathbf{P}^{(l)}(t), &
\mathbf{J}_{l,k}(t) &= \mathbf{0}\quad (k\notin\{l,l-1\}), \label{eq:dcr-J-lower}
\end{align}
for $l=2,\dots,L$. Consequently, the Jacobian has the block form
\[
\mathbf{J}(t) =
\begin{bmatrix}
\mathbf{R}^{(1)}(t) & \mathbf{0} & \cdots & \mathbf{0} & \mathbf{C}^{(1)}(t) \\
\mathbf{P}^{(2)}(t) & \mathbf{R}^{(2)}(t) & \ddots & \vdots & \mathbf{0} \\
\vdots & \ddots & \ddots & \mathbf{0} & \vdots \\
\mathbf{0} & \cdots & \mathbf{P}^{(L)}(t) & \mathbf{R}^{(L)}(t) & \mathbf{0}
\end{bmatrix},
\]
which is block lower-bidiagonal with an additional non-zero block in the top-right corner coming from the cycle.
Here $\mathbf{R}^{(l)}(t)$ are the \textit{recurrent} diagonal blocks, $\mathbf{P}^{(l)}(t)$ are the \textit{cyclic} sub-diagonal blocks and
$\mathbf{C}^{(1)}(t)$ is the cyclic top-right block.

\begin{align}
\mathbf{R}^{(l)} &= (1-\alpha)\mathbf{I}_{n_l} + \alpha \mathbf{W}_{\text{rec}}^{(l)},\\
\mathbf{P}^{(l)} &= \alpha \mathbf{W}_{\text{cyc}}^{(l-1)},\qquad l=2,\dots,L,\\
\mathbf{C}^{(1)} &= \alpha \mathbf{W}_{\text{cyc}}^{(L)}.
\end{align}

For compactness, we define the global   \DCR\@ weight matrix $\mathbf{W}_{\text{ \DCR}}\in\mathbb{R}^{N\times N}$ as
\begin{equation}
\mathbf{W}_{\text{ \DCR}} \coloneqq \WtotCycle
\label{eq:dcr-WDCR}
\end{equation}
where the off-diagonal blocks implement the cyclic couplings and the diagonal blocks are the intra-module recurrent weights.

Assuming a null input sequence $\mathbf{u}(t) = \mathbf{0}$ and analyzing the system around the zero state $\mathbf{H} = \mathbf{0}$, we have $\tanh(0)=0$ and $\tanh'(0)=1$, implying $\mathbf{D}^{(l)}(t) = \mathbf{I}_{n_l}$ at $\mathbf{H}=\mathbf{0}$.
Substituting this into the block definitions \eqref{eq:dcr-R}--\eqref{eq:dcr-C} yields


With this notation, the Jacobian of the linearized system at $\mathbf{H}=\mathbf{0}$ can be written as
\begin{equation}
\mathbf{J}_{F,\mathbf{0}} = (1-\alpha)\mathbf{I} + \alpha\,\mathbf{W}_{\text{ \DCR}}.
\label{eq:dcr-J-linearised}
\end{equation}

\section{Special case: no cycle}
\label{sec:dcr-nocycle}

If $\mathbf{W}_{\text{cyc}}^{(L)}=\mathbf{0}$ (and, more generally, $\mathbf{W}_{\text{cyc}}^{(l)}=\mathbf{0}$ for all $l$), the top-right block vanishes and $\mathbf{J}(t)$ becomes block lower-triangular. In this case the eigenvalues of $\mathbf{J}_{F,\mathbf{0}}$ are simply the union of the spectra of the diagonal blocks. The ESP then reduces to the classical layer-wise requirement
\[
\max_{l=1,\dots,L} \rho\!\big((1-\alpha)\mathbf{I} + \alpha\mathbf{W}_{\text{rec}}^{(l)}\big) < 1,
\]

\section{Special Case: no recurrent}
Instead of removing the cycle, if we remove the recurrent connections, i.e., $\mathbf{W}_{\text{rec}}^{(l)}=\mathbf{0}$ for all $l$, then the Jacobian becomes
\[
\mathbf{J}(t) = (1-\alpha)\mathbf{I} + \alpha
\begin{bmatrix} 
\mathbf{0} & \mathbf{0} & \cdots & \mathbf{0} & \mathbf{W}_{\text{cyc}}^{(L)} \\
\mathbf{W}_{\text{cyc}}^{(1)} & \mathbf{0} & \ddots & \vdots & \mathbf{0} \\
\vdots & \ddots & \ddots & \mathbf{0} & \vdots \\
\mathbf{0} & \cdots & \mathbf{W}_{\text{cyc}}^{(L-1)} & \mathbf{0} & \mathbf{0}
\end{bmatrix}.
\]

The ESP condition then depends solely on the cyclic coupling weights. This configuration in the case of single units is the same as in the paper \cite{mincomplex}, 
where we find the find same behaviour of SCR (Simple Cycle Reservoirs) on MC task. 
 a ring of interconnected feedforward layers, where each layer's output feeds into the next, forming a closed loop. The stability analysis in this case focuses on the spectral properties of the cyclic weight matrices and their influence on the overall system dynamics.

%% =========================================================================
% DEEP ANTISYMMETRIC RESERVOIR (\DAR) APPENDIX
% =========================================================================

\chapter{Jacobian Derivation for Antisymmetric modularity}
\label{ch:appendix-B}

This appendix derives the Jacobian and stability conditions for the Echo State Property (ESP) of the \emph{Deep Antisymmetric Reservoir} (DAR) described in Chapter~\ref{chap:methodology}. The derivation follows the same approach used for the  \DCR in Appendix~\ref{ch:appendix-A}.

\section{Model definition}
\label{sec:dar-definition}

Let the input be $\mathbf{u}(t)\in\mathbb{R}^{N_U}$ and consider $L$ reservoir modules with states
\[
\mathbf{h}^{(l)}(t)\in\mathbb{R}^{n_l},\qquad l=1,\dots,L.
\]
The global reservoir state is the concatenation
\[
\mathbf{H}(t) \coloneqq \big[\mathbf{h}^{(1)}(t),\dots,\mathbf{h}^{(L)}(t)\big] \in \mathbb{R}^{N},\quad N=\sum_{l=1}^L n_l.
\]

All modules share the same leak rate $\alpha\in(0,1]$, and the element-wise nonlinearity is $\tanh(\cdot)$.
As introduced in Eq.~\eqref{eq:method:antisym_leak}, the pre-activation of each module includes:
\begin{itemize}
    \item A \textbf{current input term} from the previous layer at time $t$;
    \item A \textbf{local memory term} with dissipation $-\gamma\mathbf{I}$;
    \item An \textbf{antisymmetric coupling term} with strength $\epsilon$ connecting adjacent modules at time $t-1$.
\end{itemize}

The pre-activation for the first module is:
\begin{equation}
\mathbf{z}^{(1)}(t) =
\mathbf{W}_{\text{in}}^{(1)} \mathbf{u}(t)
+ \big(\mathbf{W}_{\text{rec}}^{(1)} - \gamma \mathbf{I}\big) \mathbf{h}^{(1)}(t-1)
- \epsilon \mathbf{C}^{\top(1)} \mathbf{h}^{(2)}(t-1)
+ \mathbf{b}^{(1)}.
\label{eq:dar-z1}
\end{equation}

For intermediate modules $l=2,\dots,L-1$:
\begin{equation}
\mathbf{z}^{(l)}(t) =
\mathbf{W}_{\text{in}}^{(l)} \mathbf{h}^{(l-1)}(t)
+ \big(\mathbf{W}_{\text{rec}}^{(l)} - \gamma \mathbf{I}\big) \mathbf{h}^{(l)}(t-1)
+ \epsilon \mathbf{C}^{(l)} \mathbf{h}^{(l-1)}(t-1)
- \epsilon \mathbf{C}^{\top(l)} \mathbf{h}^{(l+1)}(t-1)
+ \mathbf{b}^{(l)}.
\label{eq:dar-zl}
\end{equation}

For the last module $l=L$:
\begin{equation}
\mathbf{z}^{(L)}(t) =
\mathbf{W}_{\text{in}}^{(L)} \mathbf{h}^{(L-1)}(t)
+ \big(\mathbf{W}_{\text{rec}}^{(L)} - \gamma \mathbf{I}\big) \mathbf{h}^{(L)}(t-1)
+ \epsilon \mathbf{C}^{(L)} \mathbf{h}^{(L-1)}(t-1)
+ \mathbf{b}^{(L)}.
\label{eq:dar-zL}
\end{equation}

Each module is updated through the leaky integration rule:
\begin{equation}
\mathbf{h}^{(l)}(t) = (1-\alpha)\,\mathbf{h}^{(l)}(t-1) + \alpha\,\tanh\big(\mathbf{z}^{(l)}(t)\big).
\label{eq:dar-leak}
\end{equation}

\section{Jacobian of the global state map of \DAR}
\label{sec:dar-jacobian}

For a fixed input $\mathbf{u}(t)$, define the state transition map
\[
\mathbf{H}(t) = F\big(\mathbf{u}(t),\mathbf{H}(t-1)\big).
\]
The Jacobian with respect to the previous global state is
\[
\mathbf{J}(t) \coloneqq \frac{\partial \mathbf{H}(t)}{\partial \mathbf{H}(t-1)} \in \mathbb{R}^{N\times N}.
\]

\subsection{Layerwise derivatives}

Let
\[
\mathbf{D}^{(l)}(t) = \mathrm{diag}\!\big(\tanh'(\mathbf{z}^{(l)}(t))\big) \in \mathbb{R}^{n_l\times n_l}.
\]

From \eqref{eq:dar-leak}--\eqref{eq:dar-zL}, we compute the partial derivatives for each module.

\paragraph{First module ($l=1$):}
\begin{align}
\frac{\partial \mathbf{h}^{(1)}(t)}{\partial \mathbf{h}^{(1)}(t-1)}
&= (1-\alpha) \mathbf{I}_{n_1} + \alpha \mathbf{D}^{(1)}(t)\, \big(\mathbf{W}_{\text{rec}}^{(1)} - \gamma \mathbf{I}\big), \label{eq:dar-j11}\\
\frac{\partial \mathbf{h}^{(1)}(t)}{\partial \mathbf{h}^{(2)}(t-1)}
&= -\alpha \epsilon \mathbf{D}^{(1)}(t)\, \mathbf{C}^{\top(1)}. \label{eq:dar-j12}
\end{align}

\paragraph{Intermediate modules ($l=2,\dots,L-1$):}
\begin{align}
\frac{\partial \mathbf{h}^{(l)}(t)}{\partial \mathbf{h}^{(l-1)}(t-1)}
&= \alpha \epsilon \mathbf{D}^{(l)}(t)\, \mathbf{C}^{(l)}, \label{eq:dar-jl-lminus1}\\
\frac{\partial \mathbf{h}^{(l)}(t)}{\partial \mathbf{h}^{(l)}(t-1)}
&= (1-\alpha) \mathbf{I}_{n_l} + \alpha \mathbf{D}^{(l)}(t)\, \big(\mathbf{W}_{\text{rec}}^{(l)} - \gamma \mathbf{I}\big), \label{eq:dar-jll}\\
\frac{\partial \mathbf{h}^{(l)}(t)}{\partial \mathbf{h}^{(l+1)}(t-1)}
&= -\alpha \epsilon \mathbf{D}^{(l)}(t)\, \mathbf{C}^{\top(l)}. \label{eq:dar-jl-lplus1}
\end{align}

\paragraph{Last module ($l=L$):}
\begin{align}
\frac{\partial \mathbf{h}^{(L)}(t)}{\partial \mathbf{h}^{(L-1)}(t-1)}
&= \alpha \epsilon \mathbf{D}^{(L)}(t)\, \mathbf{C}^{(L)}, \label{eq:dar-jL-Lminus1}\\
\frac{\partial \mathbf{h}^{(L)}(t)}{\partial \mathbf{h}^{(L)}(t-1)}
&= (1-\alpha) \mathbf{I}_{n_L} + \alpha \mathbf{D}^{(L)}(t)\, \big(\mathbf{W}_{\text{rec}}^{(L)} - \gamma \mathbf{I}\big). \label{eq:dar-jLL}
\end{align}

\subsection{Block structure}

Introduce the block matrices for convenience:
\begin{align}
\mathbf{R}^{(l)}(t) &\coloneqq (1-\alpha) \mathbf{I}_{n_l} + \alpha \mathbf{D}^{(l)}(t)\, \big(\mathbf{W}_{\text{rec}}^{(l)} - \gamma \mathbf{I}\big), \label{eq:dar-R}\\
\mathbf{F}^{(l)}(t) &\coloneqq \alpha \epsilon \mathbf{D}^{(l)}(t)\, \mathbf{C}^{(l)},\qquad l=2,\dots,L, \label{eq:dar-F}\\
\mathbf{B}^{(l)}(t) &\coloneqq -\alpha \epsilon \mathbf{D}^{(l)}(t)\, \mathbf{C}^{\top(l)},\qquad l=1,\dots,L-1. \label{eq:dar-B}
\end{align}

Here $\mathbf{R}^{(l)}(t)$ represents the diagonal (self-recurrence) blocks, $\mathbf{F}^{(l)}(t)$ the forward coupling from layer $l-1$ to layer $l$, and $\mathbf{B}^{(l)}(t)$ the backward coupling from layer $l+1$ to layer $l$.

Let $\mathbf{J}_{l,k}(t)\coloneqq \partial \mathbf{h}^{(l)}(t)/\partial \mathbf{h}^{(k)}(t-1)$. The non-zero blocks are:
\begin{align}
\mathbf{J}_{l,l}(t) &= \mathbf{R}^{(l)}(t), \qquad l=1,\dots,L, \label{eq:dar-J-diag}\\
\mathbf{J}_{l,l-1}(t) &= \mathbf{F}^{(l)}(t), \qquad l=2,\dots,L, \label{eq:dar-J-sub}\\
\mathbf{J}_{l,l+1}(t) &= \mathbf{B}^{(l)}(t), \qquad l=1,\dots,L-1. \label{eq:dar-J-super}
\end{align}

Consequently, the Jacobian has a \textbf{block tridiagonal} structure:
\[
\mathbf{J}(t) =
\begin{bmatrix}
\mathbf{R}^{(1)}(t) & \mathbf{B}^{(1)}(t) & \mathbf{0} & \cdots & \mathbf{0} \\
\mathbf{F}^{(2)}(t) & \mathbf{R}^{(2)}(t) & \mathbf{B}^{(2)}(t) & \ddots & \vdots \\
\mathbf{0} & \mathbf{F}^{(3)}(t) & \ddots & \ddots & \mathbf{0} \\
\vdots & \ddots & \ddots & \mathbf{R}^{(L-1)}(t) & \mathbf{B}^{(L-1)}(t) \\
\mathbf{0} & \cdots & \mathbf{0} & \mathbf{F}^{(L)}(t) & \mathbf{R}^{(L)}(t)
\end{bmatrix}.
\]


\subsection{Linearization at the zero state}

Assuming a null input sequence $\mathbf{u}(t) = \mathbf{0}$ and analyzing the system around the zero state $\mathbf{H} = \mathbf{0}$, we have $\tanh(0)=0$ and $\tanh'(0)=1$, implying $\mathbf{D}^{(l)}(t) = \mathbf{I}_{n_l}$ at $\mathbf{H}=\mathbf{0}$.

Substituting this into the block definitions \eqref{eq:dar-R}--\eqref{eq:dar-B} yields:
\begin{align}
\mathbf{R}^{(l)} &= (1-\alpha)\mathbf{I}_{n_l} + \alpha \big(\mathbf{W}_{\text{rec}}^{(l)} - \gamma \mathbf{I}\big), \label{eq:dar-R-lin}\\
\mathbf{F}^{(l)} &= \alpha \epsilon \mathbf{C}^{(l)},\qquad l=2,\dots,L, \label{eq:dar-F-lin}\\
\mathbf{B}^{(l)} &= -\alpha \epsilon \mathbf{C}^{\top(l)},\qquad l=1,\dots,L-1. \label{eq:dar-B-lin}
\end{align}

For compactness, we define the global \DAR\@ weight matrix $\mathbf{W}_{\text{\DAR}}\in\mathbb{R}^{N\times N}$ as:
\begin{equation}
\mathbf{W}_{\text{\DAR}} \coloneqq
\begin{bmatrix}
\mathbf{W}_{\text{rec}}^{(1)}-\gamma \mathbf{I} & -\epsilon\mathbf{C}^{\top(1)} & \cdots & \mathbf{0} \\
\epsilon\mathbf{C}^{(2)} & \mathbf{W}_{\text{rec}}^{(2)}-\gamma \mathbf{I} & \ddots & \vdots \\
\vdots & \ddots & \ddots & -\epsilon\mathbf{C}^{\top(L-1)} \\
\mathbf{0} & \cdots & \epsilon\mathbf{C}^{(L)} & \mathbf{W}_{\text{rec}}^{(L)}-\gamma \mathbf{I}
\end{bmatrix},
\label{eq:dar-WDAR}
\end{equation}
where the diagonal blocks contain the damped recurrent weights and the off-diagonal blocks implement the antisymmetric inter-module couplings.

With this notation, the Jacobian of the linearized system at $\mathbf{H}=\mathbf{0}$ can be written as:
\begin{equation}
\mathbf{J}_{F,\mathbf{0}} = (1-\alpha)\mathbf{I} + \alpha\,\mathbf{W}_{\text{\DAR}}.
\label{eq:dar-J-linearised}
\end{equation}

% =========================================================================
% Supporting Plots APPENDIX
% =========================================================================

\chapter{Supporting Materials}
\label{ch:appendix-B}
\todo{Posso aggiungere anche I plot delle curve MC degli altri modelli che non ho messo nella parte di memory capacity}
This appendix contains some plots that we did not showed in the main work, but that help us to explain our ideas about how we decided on the architecture of our models.

\section{Time Scale Differentiation}
For demonstrate the time scale differentiation property of ours models, 
with respect to standard DeepESN, we report here plots of the distance between the state of each layer, we follow the same procedure
described by \cite{GALLICCHIO201787}, where we take Artificial dataset, consisting in a one-of-10 econding, in which the input element is drawn 
from a uniform distribution from an alphabet of 10 symbols. Each symbol is represented by a 10-dimensional binary vector with a single high bit (1) and nine low bits (0).
We take two sequence, $S1$ and $S2$, that are identical except for a single symbol at $t=100$, we denote 
$h_u^{(l)}(t)$ and $h_p^{(l)}(t)$ the states of layer $l$ at time $t$ for the unperturbed sequence $S1$ and the perturbed sequence $S2$, respectively.
For each layer we define the distance between the two states as:
\[
D^{(l)}(t) = \| h_u^{(l)}(t) - h_p^{(l)}(t) \|_2
\]
In Figure \ref{fig:timescale_differentiation} we report the distance $D^{(l)}(t)$ 
for each layer $l$ of a DeepESN and of a  \DCR, we can also observe that  modules' euclidean distances shrink a little bit slower than DeepESN.
 
We can see how in the DeepESN \ref{fig:deepESN_timescale_differentiation} the perturbation effect on layers.
     The   \DCR\@ \ref{fig:dcr_timescale_differentiation} exhibits more similar states across layers after perturbation, because it keep receiving at each layer new information, plus the cyclic connection, that make some
     layers a higher distance. The \DAR\@ \ref{fig:dar_timescale_differentiation} on the other hand, having information flowing in both direction and hierarchical structure for the input, is able 
     get a more clear distinction between layers.
    \label{fig:timescale_differentiation}
\begin{figure}[h]
    \centering
    \begin{subfigure}{0.95\textwidth}
        \centering
        \includegraphics[width=\textwidth]{images/distance_states/distance_reservoir_states_DeepReservoir_5_layers.png}
        \caption{DeepESN}
        \label{fig:deepESN_timescale_differentiation}
    \end{subfigure}
    \hfill
    \begin{subfigure}{0.95\textwidth}
        \centering
        \includegraphics[width=\textwidth]{images/distance_states/distance_reservoir_states_Cycle DeepReservoir_5_layers.png}
        \caption{ \DCR}
        \label{fig:dcr_timescale_differentiation}
    \end{subfigure}
    \hfill
    % add antisymmetric
    \begin{subfigure}{0.95\textwidth}
        \centering
        \includegraphics[width=\textwidth]{images/distance_states/distance_reservoir_states_Antisymmetric DeepReservoir_5_layers.png}
        \caption{\DAR}
        \label{fig:dar_timescale_differentiation}
    \end{subfigure}
    \caption{Time scale differentiation analysis for DeepESN and  \DCR. 
    The plots show the distance $D^{(l)}(t)$ between the states of each layer $l$ for two sequences differing by a single symbol at $t=100$.}
\end{figure}

\section{Memory Capacity Additional Plots}
We report here additional plot of MC, with the same methodology we described in sec \ref{sec:exp:mc}.
From our reasoning that SCR is a special case of our more general  \DCR, we expected that a modular structure with only cyclic
connections would still perform well. Moreover by removing diagonal recurrence, we get a block-circulant matrix, which has some theoretical
advantages on spectral properties, in fact these matrices admit a spectral decomposition via the discrete Fourier transform across layers.
So the eigenvalues of it's Jacobian are given by $\lambda_{k,j} = \mu_j \omega_k$, 
where $\mu_j$ are the eigenvalues of the block and $\omega_k = e^{2\pi i k/L}$ are the $L$-th roots of unity.

% show eigenvalues
\begin{comment}
\begin{figure}[ht!]
    \centering
    \includegraphics[width=\textwidth]{images/todo.png}
    \caption{ \DCR\@ without recurrence Eigenvalues}
        \label{fig:dcr_no_recurrence_eigenvalues}
    \caption{Eigenvalues of the Jacobian of a  \DCR\@ without recurrent connections,
     with 5 layers and 20 units per layer.
      The cyclic connections are initialized with a spectral radius of 0.9.}
\end{figure}
\end{comment}

\begin{comment}
\paragraph{Spectral properties of the pure cyclic coupling.}
Consider the block-circulant matrix $W_{\text{cyc}}$ describing the inter-layer cyclic coupling in the   \DCR\@ architecture, 
in absence of intra-layer recurrent blocks. Such a matrix is not orthogonal in general, since orthogonality would require
each inter-layer block to be orthogonal and the global permutation structure to preserve orthonormality. However,
block-circulant matrices admit a spectral decomposition via the discrete Fourier transform across layers.
In the homogeneous case, where all blocks are identical ($B_l = B$), 
the eigenvalues of $W_{\text{cyc}}$ are given by $\lambda_{k,j} = \mu_j \omega_k$, 
where $\mu_j$ are the eigenvalues of $B$ and $\omega_k = e^{2\pi i k/L}$ are the $L$-th roots of unity.
Consequently, the spectrum consists of rotated copies of the spectrum of $B$, implying $\rho(W_{\text{cyc}}) = \rho(B)$.
Therefore, eigenvalues lie on the unit circle only if $|\mu_j| = 1$ for all $j$, i.e., when the inter-layer blocks are orthogonal.
In general, the cyclic structure induces rotational symmetry in the complex plane, but the radial spread is governed by 
the spectral radius of the block matrices rather than by intrinsic orthogonality.
\end{comment}